\subsection{Matlis 对偶}

在本节中，我们假设环 A 是局部 Noether 环。

定义.—— 称 A-模 I 为 Matlis A-模，如果它是内射的，如果 \(\mathfrak{m}_{\mathrm{A}}\) 是它唯一的伴随素理想，并且 \(\kappa_{\mathrm{A}}\)-向量空间 \(\operatorname{Hom}_{\mathrm{A}}(\kappa_{\mathrm{A}}, \mathrm{I})\) 的维数为 1。

设 \(e: \kappa_{\mathrm{A}} \to I\) 为 \(\kappa_{\mathrm{A}}\) 的一个内射包（A, X, 第 20 页，定理 2）。A-模 I 是一个 Matlis 模，并且每个 Matlis A-模都同构于 I（第 2 节，命题 3 的推论）。若 A 是以 K 为分式域的离散赋值环，则 A-模 K/A 是一个 Matlis 模（第 1 节，例）。若 A 是局部 Artin 环，则 A-模 A 是 Matlis 模必须且只须 A 是 Gorenstein 环（§ 3, 第 7 节，引理 1）。

设 I 为一个 Matlis A-模。对每个整数 \(n \geqslant 0\)，用 \(I_n\) 表示 I 中被 \(\mathfrak{m}_A^n\) 零化的元素组成的子 A-模。由第 2 节命题 2，A-模 I 是诸 \(I_n\) 的并，A-模 \(I_1\) 长度为 1（即同构于 \(\kappa_A\)），并且 A-模 I 是 \(I_1\) 的一个内射包；此外，分次 gr(A)-模的典范同态

\[
\beta : \mathrm{gr} ^ {\mathrm{m} _ {\mathrm{A}}} (\mathrm{I}) \longrightarrow \operatorname{Homgr} _ {\kappa_ {\mathrm{A}}} (\mathrm{gr} (\mathrm{A}), \mathrm{I} _ {1})\tag{3}
\]

是同构。由第 2 节命题 3，I 的 A-模结构唯一地延拓为一个 \(\widehat{\mathrm{A}}\) -模结构，且典范同态 \(\widehat{\mathrm{A}} \to \operatorname{End}_{\mathrm{A}}(\mathrm{I})\) 是双射。

引理 3.—— 设 I 为一个 Matlis A-模。那么：

\begin{enumerate}
\def\labelenumi{\alph{enumi})}
\item
  I 是一个 A-hat-Matlis 模；
\item
  A-模 I 是 Artin 的，并且是一个余生成元（A, X, 第 18 页，定义 3）。
\end{enumerate}

由于 A-模 I 是内射的，\(\mathbf{A} / \mathfrak{m}_{\mathbf{A}}^{n}\)-模 \(\mathrm{I}_n\) 对每个 \(n\) 都是内射的（\(n^{\circ}2\)，引理 1，a)）。由于 \(\mathrm{I}_n\) 是 I 中被 \(\mathfrak{m}_{\widehat{\mathrm{A}}}^{n}\) 零化的元素组成的集合，\(\widehat{\mathrm{A}}\)-模 I 是内射的（引理 1，b)）。它在 \(\widehat{\mathrm{A}}\) 上不可分解，因为它在 \(\mathrm{A}\) 上已是如此；由于它包含子 \(\widehat{\mathrm{A}}\) -模 \(\mathrm{I}_1\) （同构于 \(\kappa_{\mathrm{A}}\) ），有 \(\mathfrak{m}_{\widehat{\mathrm{A}}} \in \operatorname{Ass}_{\widehat{\mathrm{A}}}(\mathrm{I})\) ，从而 \(\operatorname{Ass}_{\widehat{\mathrm{A}}}(\mathrm{I}) = \{\mathfrak{m}_{\widehat{\mathrm{A}}}\}\) （命题 1），由此得 a)。

我们现在来证明 I 是 Artin 的。对 I 的每个子 A-模 M，结合 gr(A) 的分次理想 \(\mathfrak{a}_{\mathrm{M}}\)，其定义如下：gr(A) 的元素 \(_n\) 属于 \((\mathfrak{a}_{\mathrm{M}})_{n}\)，如果它被所有线性型 \(\boldsymbol{\beta}(x)\) 零化，其中 x 取遍 \(\left((\mathrm{M}\cap\mathrm{I}_{n+1})+\mathrm{I}_{n}\right)/\mathrm{I}_{n}\)。设 M、N 为 I 的子模，使得 \(N\subset M\)；则有 \(a_{M}\subset a_{N}\)。假设 \(a_{M}=a_{N}\)；则对每个 n 有 \((M\cap I_{n+1})+I_{n}=(N\cap I_{n+1})+I_{n}\)，因为 \(\beta\) 是同构。对 n 用归纳法推出，对每个 n 有 \(M\cap I_{n+1}=N\cap I_{n+1}\)，最终得 M=N。

于是，设 \(M_{0} \supset M_{1} \supset \ldots \supset M_{n} \supset \ldots\) 为 I 的子 A-模的一个递减序列；递增序列 \(a_{M_{0}} \subset a_{M_{1}} \subset \ldots\) 是平稳的，因为 \(\text{gr}(A)\) 是一个有限型的 \(\kappa_{A}\)-代数。从而序列 \((M_{i})_{i \geqslant 0}\) 是平稳的，这蕴含 A-模 I 是 Artin 的。最后，由 A, X, 第 18 页，命题 12，A-模 I 是一个余生成元。

设 M 为 A-模。回忆（A, VIII, § 4, 第 6 节）M 的底座是 M 的单子模之和，即 M 中被 \(\mathfrak{m}_{\mathrm{A}}\) 零化的元素组成的集合；它是一个 \(\kappa_{\mathrm{A}}\)-向量空间，典范地同构于 \(\operatorname{Hom}_{\mathrm{A}}(\kappa_{\mathrm{A}}, \mathrm{M})\)。

引理 4. 设 I 为 A 上的一个 Matlis 模，设 M 为一个 \(A\)-模。下列条件等价：

\begin{enumerate}
\def\labelenumi{(\roman{enumi})}
\item
  M 是 Artin 的；
\item
  M 的每个元素都被 \(\mathfrak{m}_{\mathrm{A}}\) 的某个幂零化，且 M 的底座在 \(\kappa_{\mathrm{A}}\) 上维数有限；
\item
  存在一个整数 \(n \geqslant 0\) 和一个从 M 到 \(\mathrm{I}^n\) 中的单射 A-线性映射。
\end{enumerate}

当这些条件满足时，M 的每个内射包都同构于 \(I^{s}\)，其中 s 是 M 的底座在 \(\kappa_{A}\) 上的维数。

(iii)⇒(i)：这是显然的，因为 A-模 I 是 Artin 的（引理 3）。

(i)⇒(ii)：设 M 是 Artin 的。设 \(x \in M\)；M 的子模的递减序列 \(m_{A}^{n}x\) 是平稳的。设 n 为使 \(m_{A}^{n+1}x = m_{A}^{n}x\) 的整数；Nakayama 引理蕴含 \(m_{A}^{n}x = 0\)。此外，M 的底座作为 A-模是 Artin 的，因而作为 \(\kappa_{A}\)-向量空间也是 Artin 的，这就是说它是有限维的。

(ii)⇒(iii)：设条件 (ii) 满足；设 \(e:M\to J\) 为 M 的一个内射包。我们有 \(\operatorname{Ass}(M)\subset\{\mathfrak{m}_{A}\}\)，从而 \(\operatorname{Ass}(J)\subset\{\mathfrak{m}_{A}\}\)（\(n^{\circ} 1\)，见注 2），并且 J 同构于 I \(^{(c)}\)，其中 c 是一个基数（\(n^{\circ} 1\)，定理 1）。设 \(x\) 为 J 中被 \(\mathfrak{m}_{A}\) 零化的一个非零元素；由于 A-模 Ax 是单模且它与 \(e(M)\) 的交不退化为 0，\(x\) 属于 \(e(M)\)。于是 \(e\) 诱导一个从 M 的底座到 J 的底座上的同构；因此 M 的底座维数为 c，这就证明了 (iii) 以及最后一个断言。

引理 5.—— 每个 Artin 的 A-hat-模作为 Λ-模都是 Artin 的。

设 M 是一个 Artin 的 \(\widehat{\mathbf{A}}\)-模；M 的每个元素都被 \(\mathfrak{m}_{\widehat{\Lambda}}\) 的某个幂零化，从而被 \(\mathfrak{m}_{\mathbf{A}}\) 的某个幂零化。由第 2 节的引理 2，M 的子 A-模就是它的子 \(\widehat{\mathbf{A}}\)-模，故 M 作为 A-模是 Artin 的。

现在固定一个 Matlis A-模 I。对每个 A-模 M，用 \(D_{\mathrm{A}}(\mathrm{M})\) 表示 \(\widehat{A}\)-模

\[
\mathrm{D} _ {\mathrm{A}} (\mathrm{M}) = \operatorname{Hom} _ {\mathrm{A}} (\mathrm{M}, \mathrm{I}).
\]

把 \(\widehat{\mathrm{A}}\)-模 \(D_{\mathrm{A}}(A)\) 典范地等同于 I，把 \(\widehat{\mathrm{A}}\)-模 \(D_{\mathrm{A}}(\mathrm{I})\) 等同于 \(\widehat{\mathrm{A}}\)（\(n^{\circ}\) 2，命题 3），并把 \(\widehat{\mathrm{A}}\)-模 \(D_{\mathrm{A}}(\kappa_{\mathrm{A}})\) 等同于 \(I_{1}\)（同处）。

对任何 A-线性映射 \(f: M \rightarrow N\)，我们将用

\[
\mathrm{D} _ {\mathrm{A}} (f): \mathrm{D} _ {\mathrm{A}} (\mathrm{N}) \rightarrow \mathrm{D} _ {\mathrm{A}} (\mathrm{M})
\]

表示 \(\widehat{\mathrm{A}}\)-线性的映射 \(\operatorname{Hom}_{\Lambda}(f,1_{\mathrm{I}})\)。由于 \(\Lambda\)-模 I 是内射的，序列 \((\mathrm{D}_{\Lambda}(g),\mathrm{D}_{\mathrm{A}}(f))\) 对 A-线性映射的每个正合序列 \((f,g)\) 都正合。

\((D_{\hat{A}}(y), D_{\hat{A}}(f))\) 对映射的每个正合序列 \((f, y)\) 都正合。

我们将把这些定义应用于赋有 Matlis 模 I 的环 \(\hat{A}\)（引理 3，a)）；对每个 \(\hat{A}\)-模 P，\(D_{\hat{A}}(P)\) 因而是子 \(\hat{A}\)-模 \(\text{Hom}_{\hat{A}}(P, I)\) ，作为 \(D_{A}(P)\) 的子模。说 P 作为 A-模是 Artin 的与作为 \(\hat{A}\)-模是 Artin 的是一回事（引理 5）；在这种情形下，有 \(D_{\hat{A}}(P) = D_{A}(P)\)（同处）。

设 M 为 A-模。对 \(m \in M\)，映射 \(f \mapsto f(m)\) （属于 \(\mathrm{D}_{\mathrm{A}}(\mathrm{M})\) ，取值于 I）是 \(\widehat{\mathrm{A}}\) -线性的；用 \(\alpha_{\mathrm{M}}(m)\) 表示它。

\[
\alpha_ {\mathrm{M}}: \mathrm{M} \longrightarrow \mathrm{D} _ {\widehat {\mathrm{A}}} (\mathrm{D} _ {\mathrm{A}} (\mathrm{M})) .
\]

用 \(\widehat{\alpha}_{\mathrm{M}}:\widehat{\mathrm{A}}\otimes_{\mathrm{A}}\mathrm{M}\longrightarrow \mathrm{D}_{\widehat{\mathrm{A}}}(\mathrm{D}_{\mathrm{A}}(\mathrm{M}))\) 表示 \(\widehat{\mathrm{A}}\)-线性映射，它由 \(\alpha_{\mathrm{M}}\)

定理 2.—— 设 M 为 A-模。

\begin{enumerate}
\def\labelenumi{\alph{enumi})}
\item
  M 是 Artin 的，必须且只须 \(\hat{\mathrm{A}}\)-模 \(\mathrm{D}_{\mathrm{A}}(\mathrm{M})\) 是有限型的。在这种情形下，同态 \(\alpha_{\mathrm{M}}\) 是双射。
\item
  M 是有限型的，必须且只须 \(D_{A}(M)\) 是 Artin 的（作为 A-模或作为 \(\hat{A}\)-模）。在这种情形下同态 \(\hat{\alpha}_{M}\) 是一个同构。
\item
  M 是有限长度的，必须且只须 \(\mathrm{D}_{\mathrm{A}}(\mathrm{M})\) 是有限长度的（作为 A-模或作为 \(\widehat{A}\)-模）。在这种情形下，\(\alpha_{M}\) 是从 M 到 \(\mathrm{D}_{\mathrm{A}}(\mathrm{D}_{\mathrm{A}}(\mathrm{M}))\) 上的一个同构，并且有 \(\mathrm{long}_{\mathrm{A}}(\mathrm{D}_{\mathrm{A}}(\mathrm{M})) = \mathrm{long}_{\mathrm{A}}(\mathrm{M})\)。
\end{enumerate}

我们首先证明：对每个 A-模 M，同态 \(\alpha_{\mathrm{M}}\) 是单射。设 \(m\) 为 M 的非零元素；它的零化子含于 \(\mathfrak{m}_{\mathrm{A}}\)。因此存在一个从 \(Am\) 到 \(\kappa_{\mathrm{A}}\) 上的满射 A-同态，从而存在一个从 \(Am\) 到 I 中的非零同态。由于 I 是内射的，它延拓为一个同态 \(f:M\to I\)，使得 \(f(m)\neq0\)。这就证明了 \(\alpha_{\mathrm{M}}\) 的单射性。

设 A-模 M 是 Artin 的。由引理 4，存在一个整数 \(r\) 和一个单射 A-线性映射 \(f: \mathrm{M} \to \mathrm{I}^r\)。于是同态 \(\mathrm{D}_{\mathrm{A}}(f): \mathrm{D}_{\mathrm{A}}(\mathrm{I}^r) \longrightarrow \mathrm{D}_{\Lambda}(\mathrm{M})\) 是满射；由于 \(\mathrm{D}_{\mathrm{A}}(\mathrm{I}^r)\) 等同于 \(\widehat{\mathrm{A}}^r\) ，这就证明了 \(\widehat{\mathrm{A}}\)-模 \(\mathrm{D}_{\mathrm{A}}(\mathrm{M})\) 是有限型的。类似地，若 M 是有限型的，则存在一个整数 \(n\) 和一个满射同态 \(u: \mathrm{A}^n \to \mathrm{M}\)；同态 \(\mathrm{D}_{\Lambda}(u): \mathrm{D}_{\mathrm{A}}(\mathrm{M}) \to \mathrm{I}^n\) 是单射，因而 \(\mathrm{D}_{\mathrm{A}}(\mathrm{M})\) 是 Artin 的（作为 A-模或作为 \(\widehat{\mathrm{A}}\)-模）。

现在设 \(\widehat{\mathrm{A}}\)-模 \(\mathrm{D}_{\mathrm{A}}(\mathrm{M})\) 是 Artin 的；对典范地同构于它的 \(\widehat{\mathrm{A}}\)-模 \(\mathrm{D}_{\widehat{\mathrm{A}}}(\widehat{\mathrm{A}} \otimes_{\mathrm{A}} \mathrm{M})\) 同样如此。由前所述，\(\widehat{\mathrm{A}}\)-模 \(\mathrm{D}_{\widehat{\mathrm{A}}}(\mathrm{D}_{\widehat{\mathrm{A}}}(\widehat{\mathrm{A}} \otimes_{\mathrm{A}} \mathrm{M}))\) 是有限型的，而对 \(\widehat{\mathrm{A}} \otimes_{\mathrm{A}} \mathrm{M}\) ——它同构于 \(\mathrm{D}_{\widehat{\mathrm{A}}}(\mathrm{D}_{\widehat{\mathrm{A}}}(\widehat{\mathrm{A}} \otimes_{\mathrm{A}} \mathrm{M}))\) 的一个子模——同样如此。于是 M 是一个有限生成的 A-模（I, § 3, 第 6 节，命题 11 及 III, § 3, 第 5 节，命题 9）。类似地，若 \(\mathrm{D}_{\mathrm{A}}(\mathrm{M})\) 是一个 \(\widehat{\mathrm{A}}\)-模且是有限型的，则 \(\mathrm{D}_{\widehat{\mathrm{A}}}(\mathrm{D}_{\mathrm{A}}(\mathrm{M}))\) 是一个 \(\widehat{\mathrm{A}}\)-模。由前所述，它是一个 Artin 模，因而是 Artin 的 A-模（引理 5），对 M 同样如此。最后，有限长度的模就是有限生成的 Artin 模（A, VIII, § 1, 第 1 节，命题 1），因此 \(\mathrm{D}_{\mathrm{A}}(\mathrm{M})\) 是有限长度的，必须且只须 M 是有限长度的。

设 M 是 Artin 的。存在一个整数 r 和一个单射 A-线性映射 \(f : M \to I^{r}\)；由于 I 是 Artin 的（引理 3），A-模 Coker(f) 也是 Artin 的，于是可以找到一个整数 s 以及 A-模的一个正合序列

\[
0 \to \mathrm{M} \stackrel {f} {\longrightarrow} \mathrm{I} ^ {r} \stackrel {g} {\longrightarrow} \mathrm{I} ^ {s}.
\]

我们导出一个各行正合的交换图

\[
\begin{array}{c c c c c c} 0 \longrightarrow & \mathrm{M} & \xrightarrow {f} & \mathrm{I} ^ {r} & \xrightarrow {g} & \mathrm{I} ^ {s} \\ & \alpha_ {\mathrm{M}} \Biggl \downarrow & & \alpha_ {\mathrm{I} ^ {r}} \Biggl \downarrow & & \alpha_ {\mathrm{I} ^ {s}} \Biggl \downarrow \\ 0 \to & \mathrm{D} _ {\widehat {\mathrm{A}}} (\mathrm{D} _ {\mathrm{A}} (\mathrm{M})) & \xrightarrow {\mathrm{D} _ {\widehat {\mathrm{A}}} (\mathrm{D} _ {\Lambda} (f))} & \mathrm{D} _ {\widehat {\mathrm{A}}} (\mathrm{D} _ {\Lambda} (\mathrm{I} ^ {r})) & \xrightarrow {\mathrm{D} _ {\widehat {\mathrm{A}}} (\mathrm{D} _ {\Lambda} (g))} & \mathrm{D} _ {\widehat {\Lambda}} (\mathrm{D} _ {\Lambda} (\mathrm{I} ^ {s})) \end{array} .
\]

把 \(\widehat{\mathbf{A}}\)-模 \(\mathrm{D}_{\widehat{\mathbf{A}}}(\mathrm{D}_{\mathbf{A}}(\mathrm{I}))\) 等同于 I，并把 \(\alpha_{\mathbf{I}}\) 等同于恒同映射；从而 \(\alpha_{\mathbf{I}^r}\) 与 \(\alpha_{\mathbf{I}^s}\) 都是双射，对 \(\alpha_{\mathbf{M}}\) 同样如此（A, X, 第 7 页，推论 3）。

若 A-模 M 是有限生成的，则存在整数 m 和 n 以及 A-模的一个正合序列

\[
\mathrm{A} ^ {m} \longrightarrow \mathrm{A} ^ {n} \longrightarrow \mathrm{M} \rightarrow 0;
\]

我们导出一个各行正合的交换图：

\[
\begin{array}{c c c c c} \widehat {\mathrm{A}} ^ {m} & \longrightarrow & \widehat {\mathrm{A}} ^ {n} & \longrightarrow & \widehat {\mathrm{A}} \otimes_ {\mathrm{A}} \mathrm{M} \longrightarrow 0 \\ \widehat {\alpha} _ {\mathrm{A} ^ {m}} \Bigg \downarrow & & \widehat {\alpha} _ {\mathrm{A} ^ {n}} \Bigg \downarrow & & \widehat {\alpha} _ {\mathrm{M}} \Bigg \downarrow \\ \widehat {\mathrm{A}} ^ {m} & \longrightarrow & \widehat {\mathrm{A}} ^ {n} & \longrightarrow & \mathrm{D} _ {\mathrm{A}} (\mathrm{D} _ {\mathrm{A}} (\mathrm{M})) \to 0. \end{array}
\]

由于 \(\widehat{\alpha}_{\mathrm{A}}\) 等于 \(1_{\widehat{\mathrm{A}}}\)，由此推出 \(\widehat{\alpha}_{\mathrm{M}}\) 是同构。

剩下要证明等式 \(\operatorname{long}_{\mathrm{A}}(\mathrm{M}) = \operatorname{long}_{\mathrm{A}}(\mathrm{D}_{\Lambda}(\mathrm{M}))\) 在 M 是有限长度时成立。我们可以假设 \(M \neq 0\) ；于是存在一个正合序列

\[
0 \rightarrow \kappa_ {\mathrm{A}} \longrightarrow \mathrm{M} \longrightarrow \mathrm{N} \rightarrow 0,
\]

由此导出一个正合序列

\[
0 \rightarrow \mathrm{D} _ {\mathrm{A}} (\mathrm{N}) \longrightarrow \mathrm{D} _ {\mathrm{A}} (\mathrm{M}) \longrightarrow \mathrm{D} _ {\mathrm{A}} (\kappa_ {\mathrm{A}}) \rightarrow 0.
\]

我们有 long \(_{A}\)(M) = long \(_{A}\)(N) + 1，并且

\[
\operatorname{long} _ {\Lambda} \left(\mathrm{D} _ {\Lambda} (\mathrm{M})\right) = \operatorname{long} _ {\mathrm{A}} \left(\mathrm{D} _ {\mathrm{A}} (\mathrm{N})\right) + \operatorname{long} _ {\mathrm{A}} \left(\mathrm{D} _ {\mathrm{A}} \left(\kappa_ {\mathrm{A}}\right)\right) = \operatorname{long} _ {\mathrm{A}} \left(\mathrm{D} _ {\mathrm{A}} (\mathrm{N})\right) + 1;
\]

对整数 long \(_{A}\) (M) 用归纳法即得结论。

注.—— 设环 A 是 Artin 的。我们有 \(\text{long}_{\text{A}}(\text{I}) = \text{long}_{\text{A}}(\text{D}_{\text{A}}(A)) = \text{long}(\text{A})\)（定理 2, c)）。设 M 为一个有限生成 A-模；它有一个同构于 I \(^{s}\) ，其中 s 是 M 的底座的维数（引理 4）。因此我们有 \(\text{long}_{\text{A}}(\text{M}) \leqslant s \text{ long}(\text{A})\)；等号成立必须且只须 M 是内射的。特别地，A-模 A 是内射的，必须且只须它的底座维数为 1；这样我们就重新得到了 § 3, 第 7 节的引理 1。

